%Version 3.1 December 2024 -- Springer Nature sn-jnl.cls
\documentclass[pdflatex,sn-basic,Numbered]{sn-jnl}

\usepackage{graphicx}
\usepackage{multirow}
\usepackage{amsmath,amssymb,amsfonts}
\usepackage{amsthm}
\usepackage{mathrsfs}
\usepackage[title]{appendix}
\usepackage{xcolor}
\usepackage{textcomp}
\usepackage{manyfoot}
\usepackage{algorithm}
\usepackage{algorithmicx}
\usepackage{algpseudocode}
\usepackage{listings}
\usepackage{enumitem}
\setlist[itemize]{leftmargin=1.5em,itemsep=0.15em,topsep=0.15em}
\setlist[enumerate]{leftmargin=1.5em,itemsep=0.15em,topsep=0.15em}

\theoremstyle{thmstyleone}
\newtheorem{theorem}{Theorem}
\newtheorem{proposition}[theorem]{Proposition}
\theoremstyle{thmstyletwo}
\newtheorem{remark}{Remark}

\raggedbottom
\setlength{\textfloatsep}{8pt plus 2pt minus 2pt}
\setlength{\intextsep}{6pt plus 2pt minus 2pt}
\setlength{\floatsep}{6pt plus 2pt minus 2pt}
\setlength{\abovecaptionskip}{4pt}
\setlength{\belowcaptionskip}{2pt}
\setlength{\abovedisplayskip}{6pt plus 2pt minus 2pt}
\setlength{\belowdisplayskip}{6pt plus 2pt minus 2pt}
\setlength{\abovedisplayshortskip}{3pt plus 2pt}
\setlength{\belowdisplayshortskip}{3pt plus 2pt}

\begin{document}

\title[Inverse PINNs for Synaptic Dopamine Transport]{When Inverse
Physics-Informed Neural Networks Become Practically Identifiable:
A Case Study in Synaptic Dopamine Transport}

\author*[1]{\fnm{Emmanuel} \sur{Hackman}}\email{Emmanuel.Hackman@usm.edu}

\author[1]{\fnm{Huiqing} \sur{Zhu}}\email{Huiqing.Zhu@usm.edu}

\affil*[1]{\orgdiv{School of Mathematics and Natural Sciences}, \orgname{University of Southern Mississippi}, \orgaddress{\city{Hattiesburg}, \postcode{39406}, \state{MS}, \country{USA}}}

\abstract{Quantifying dopamine clearance matters for understanding
dopaminergic dysfunction in Parkinson's disease, but classical
mesh-based methods struggles under sparse, noisy observation. We formulate recovery of
the diffusion coefficient $D$ and reuptake rate $k$ as a
parameter inverse problem for a 2D reaction-diffusion PDE of
synaptic dopamine transport, using physics-informed neural
networks (PINNs) as a mesh-free forward and inverse solver. All
validation uses synthetic data from a finite-difference
reference; no biological measurements are used yet. The forward
model attains an L2 relative error of $1.52\%$ against the
finite-difference reference and $3.75\%$ against the analytical
solution. For the inverse problem, evaluated over a five-seed
ensemble, we recover $D = 0.3515 \pm 0.0034$
$\mu\text{m}^2/\text{ms}$ and $k = 0.0573 \pm 0.0006$
$\text{ms}^{-1}$, with mean relative errors of $9.85\%$ and
$14.61\%$. The inverse problem is well-conditioned along one
direction and ill-conditioned along a diffusion-reuptake trade
ray, with implications for real-data experimental design. Unlike
a Levenberg-Marquardt baseline, which drives $k$ to non-physical
negative values under bounded-domain data, the PINN's
residual-based formulation recovers physically meaningful
parameters. This identifiability-aware framework is a reusable
tool for inverse-PDE recovery of synaptic neurotransmitter
dynamics; application to experimental fast-scan cyclic
voltammetry data is the immediate next step.}

\keywords{Physics-informed neural networks, dopamine diffusion,
reaction-diffusion equations, inverse problem, Parkinson's
disease}

\maketitle

\noindent\textbf{Manuscript counts (per journal Title Page
requirements):} Total word count (title page through Conclusion,
excluding Declarations and reference list): approximately
5{,}800. Abstract word count: 198. Number of figures: 8. Number
of tables: 6. Corresponding author telephone: +1 (601) 266-4934
(School of Mathematics and Natural Sciences office); fax: not
applicable.

\noindent\textbf{Author biographies} \\[0.5em]
Emmanuel Hackman is a PhD student in Mathematics at the
University of Southern Mississippi, researching physics-informed
Machine Learning for differential equations. \\[0.5em]
Huiqing Zhu is Professor of Mathematics at the University of
Southern Mississippi,
specializing in numerical methods for differential equations.

\section*{Glossary of Terms}
\label{sec:glossary}

\begin{tabular}{@{}lll@{}}
$C(x,y,t)$ & dopamine concentration field & $\mu$M \\
$D$        & effective diffusion coefficient & $\mu\text{m}^2/\text{ms}$ \\
$k$        & linearised dopamine transporter(DAT) reuptake rate constant & $\text{ms}^{-1}$ \\
$L$        & domain side length & $\mu$m \\
$T$        & simulation time window & ms \\
$\sigma$   & release-pulse half-width & $\mu$m \\
$C_0$      & peak release concentration & $\mu$M \\
\end{tabular}

\section{Introduction}
\label{sec:intro}

Parkinson's disease (PD) is the second most prevalent
neurodegenerative disorder worldwide, affecting approximately
10 million people and imposing a substantial public-health
burden \cite{GBD2016}. Its hallmark pathology is progressive
degeneration of dopaminergic neurons in the substantia nigra
pars compacta, reducing dopamine availability in the striatum
and producing the cardinal motor symptoms of PD: bradykinesia,
rigidity, and tremor \cite{Hornykiewicz1966, DeLong2007}.

Dopaminergic signaling begins when an action potential triggers
$\text{Ca}^{2+}$-dependent vesicle fusion, releasing
$\sim$$10^3$--$10^4$ dopamine molecules into the $\sim$20\,nm
synaptic cleft~\cite{Wiencke2020}, which then diffuses with
effective coefficient $D \approx 0.3\,\mu\text{m}^2/\text{ms}$
\cite{Nicholson1981, CraggRice2004} and is cleared by the
dopamine transporter (DAT), modeled at sub-saturating
concentrations as a first-order reuptake rate
$k$~\cite{CraggRice2004}. PD reduces both functional release-site
density and striatal DAT expression~\cite{Rice2008,
CraggRice2004}, so changes in $(D, k)$ are a quantitatively
meaningful proxy for disease state. Fast-scan cyclic voltammetry
(FSCV) provides millisecond-resolution, sub-micrometer dopamine
recordings \emph{in vivo} and \emph{ex vivo}~\cite{Rodeberg2017},
but yields exactly the sparse, noisy regime in which recovering
$(D, k)$ is hardest.

Classical reaction-diffusion models of dopamine transport
\cite{Nicholson1981, Rice2008, Wiencke2020} use finite-difference
or finite-element solvers, requiring costly mesh construction in
complex neuroanatomical geometries and repeated forward solves
per optimization step, an approach that scales poorly and is
ill-suited to sparse experimental data. Physics-informed neural
networks (PINNs)~\cite{Raissi2019} instead embed the governing
PDE directly into the network's training objective, solving
forward and inverse problems mesh-free. PINNs have been applied
to cardiac electrophysiology~\cite{Sahli2020}, systems-biology
kinetics~\cite{Yazdani2020}, and PDE discovery from neural
imaging~\cite{Rudy2017}, but, to our knowledge and conditional
on the limits of our search, not to synaptic neurotransmitter
diffusion.

This paper makes four contributions. First, an inverse-PINN
formulation for spatially-resolved 2D dopamine transport,
extending prior inverse-PINN work in systems
biology~\cite{Yazdani2020}, lumped ODE kinetics without spatial
structure to a spatially explicit problem. Second, forward
validation against both an analytical solution and a
finite-difference reference, with relative $L^2$ errors
comparable to established numerical schemes without a spatial
mesh. Third, the inverse problem recovering $D$ and $k$ from
sparse, noisy synthetic data, characterizing its practical
identifiability: well-conditioned along one direction,
ill-conditioned along a diffusion-reuptake trade ray, with
direct implications for experimental design; we further show the
PINN's bounded-domain residual formulation succeeds where a
classical Levenberg-Marquardt baseline fails, driving $k$
negative under the same data. Fourth, we release the complete
implementation seed-controlled scripts and a reproducible
notebook as an open-source, ready-to-extend baseline for
application to experimental FSCV data, the immediate next step
toward a patient-specific tool for dopaminergic dysfunction in
PD.

\section{Methods}
\label{sec:methods}

\subsection{Model Formulation}
\label{sec:math}

We model dopamine concentration $C(x,y,t)$ on a 2D domain
spanning the synaptic cleft and surrounding perisynaptic space.
The 2D restriction is motivated by the cleft's thinness
($h \approx 20$--$30\,\text{nm}$, two orders of magnitude below
the lateral spreading scale $\sqrt{2DT} \approx 3.6\,\mu\text{m}$,
so vertical equilibration is essentially instantaneous), by
consistency with prior striatal dopamine
modeling~\cite{Wiencke2020}, and because it admits a closed-form
analytical reference. Extension to 3D is a
straightforward modification of the input layer and PDE residual
(Section~\ref{sec:disc_future}).

Following Cragg and Rice~\cite{CraggRice2004}, DAT action at
sub-saturating concentrations is approximated by a first-order
linear sink $-k\,C$ (accurate when $C \ll K_m \approx 0.2\,\mu$M),
reducing recovery to two scalars $(D,k)$. Release is modeled as a
radially symmetric Gaussian pulse at $t=0$, consistent with the
sub-millisecond timescale of vesicular
exocytosis~\cite{Wiencke2020}.

With these modeling choices, we consider the reaction-diffusion
initial-boundary-value problem on the square domain
$\Omega = (-L/2, L/2)^2$, $t \in (0, T]$,

\begin{equation}
  \frac{\partial C}{\partial t} = D\,\nabla^2 C \;-\; k\,C,
  \quad (x, y) \in \Omega,\ t \in (0, T],
  \label{eq:pde}
\end{equation}
where $\nabla^2 C = \partial_{xx} C + \partial_{yy} C$ is the
two-dimensional Laplacian, with the radially symmetric Gaussian
initial condition

\begin{equation}
  C(x, y, 0) = C_0 \exp\!\left( -\frac{x^2 + y^2}{2\sigma^2} \right),
  \quad (x, y) \in \overline{\Omega},
  \label{eq:ic}
\end{equation}

and zero-flux (Neumann) boundary conditions on all four edges
of $\partial \Omega$,
\begin{align}
  \left.\frac{\partial C}{\partial x}\right|_{x=\pm L/2}
    &= 0, \quad y \in [-L/2, L/2],\ t \in (0, T],
    \label{eq:bc_left} \\
  \left.\frac{\partial C}{\partial y}\right|_{y=\pm L/2}
    &= 0, \quad x \in [-L/2, L/2],\ t \in (0, T].
    \label{eq:bc_right}
\end{align}

The zero-flux conditions \eqref{eq:bc_left}--\eqref{eq:bc_right}
encode chemical isolation from neighboring synaptic compartments
on the timescales of interest. The Gaussian profile in
\eqref{eq:ic} represents the instantaneous footprint of a single
quantal release event, with $\sigma$ the release radius and
$C_0$ the peak concentration.

All model parameters are taken from the experimental dopamine
literature for striatal tissue~\cite{CraggRice2004, Nicholson1981}
and are summarized in Table~\ref{tab:params}.

\begin{table}[t]
\centering
\caption{\textbf{Model parameters for the dopamine synaptic
diffusion problem.} $D$, $k$ from canonical experimental values
for striatal tissue; $L$, $\sigma$, $C_0$ from the Wiencke et
al.~\cite{Wiencke2020} model; $T$ is a modelling choice.}
\label{tab:params}
\begin{tabular}{llll}
\toprule
\textbf{Parameter} & \textbf{Symbol} & \textbf{Value} & \textbf{Source} \\
\midrule
Diffusion coefficient    & $D$   & $0.32$ $\mu\text{m}^2/\text{ms}$ & \cite{Nicholson1981, CraggRice2004} \\
Reuptake rate constant   & $k$   & $0.05$ $\text{ms}^{-1}$           & \cite{CraggRice2004} \\
Domain length            & $L$   & $5.0$ $\mu\text{m}$               & \cite{Wiencke2020} \\
Simulation time          & $T$   & $20$ ms                            & (model choice) \\
Peak source amplitude    & $C_0$ & $1.0$ $\mu$M                       & \cite{Wiencke2020} \\
\botrule
\end{tabular}
\end{table}

\subsection{Implementation and Reproducibility}
\label{sec:implementation}

Detailed JAX/Flax NNX code listings for the network module, PDE
residual, and inverse-problem training loop are provided in
Online Resource 1; the closed-form analytical solution
derivation underlying the forward-problem validation reference
is provided in Online Resource 2.

\begin{figure}[t]
  \centering
  \includegraphics[width=0.88\textwidth]{02_Figures/Fig1_preview.png}
  \caption{\textbf{Overview of the inverse-PDE problem and the
  physics-informed neural network methodology.}
  \textbf{A.} Forward setting: 2D reaction-diffusion PDE on
  $\Omega = [-L/2, L/2]^2$ with Neumann boundaries and a Gaussian
  release pulse; crosses mark a subset of the
  $N_{\mathrm{obs}} = 400$ sparse noisy observations.
  \textbf{B.} Inverse PINN: a Flax NNX MLP
  $\mathcal{N}_\theta : (x, y, t) \mapsto \hat{C}$ is trained to
  satisfy the PDE residual while fitting sparse noisy
  observations; $D$ and $k$ are log-parametrized learnable
  scalars optimized jointly with the network weights.
  \textbf{C.} Identifiability: a one-dimensional ``trade ray''
  through $(D, k)$ space along which the loss is flat to within
  the noise scale. Five seeds (red dots) cluster along this ridge
  above ground truth $(D^*, k^*) = (0.32, 0.05)$, reflecting the
  structural limit of Section~\ref{sec:identifiability}}
  \label{fig:schematic}
\end{figure}

\subsection{PINN Architecture}
\label{sec:arch}

We employ a fully-connected feedforward network
$\mathcal{N}_\theta : \mathbb{R}^3 \to \mathbb{R}$ mapping
space-time coordinates $(x,y,t)$ to predicted concentration
$\hat{C}(x,y,t;\theta)$, with four hidden layers of $64$ neurons
each and $\tanh$ activation (continuously differentiable, the
standard PINN choice~\cite{Raissi2019}):
\begin{equation}
  \hat{C} = \mathcal{N}_\theta(x, y, t) =
  (W_\ell \circ \sigma \circ W_{\ell-1} \circ \cdots
  \circ \sigma \circ W_1)(x, y, t),
  \label{eq:network}
\end{equation}
with $\ell = 5$ weight matrices (four hidden layers plus a
linear output layer, allowing $\hat{C}$ unbounded range). Weights
use Glorot initialization~\cite{Glorot2010}; the network has
$12{,}801$ trainable parameters in total.

Spatial and temporal derivatives in the PDE residual
\eqref{eq:pde} are evaluated by automatic differentiation of
$\mathcal{N}_\theta$, removing any need for finite-difference
stencils. The implementation uses JAX~\cite{JAX2018} with the
network as a Flax NNX module~\cite{Heek2020}, Adam via
Optax~\cite{Optax2020}, and L-BFGS via
jaxopt~\cite{Blondel2022}; JIT compilation gives a
2--5$\times$ speed-up over a comparable PyTorch implementation.
Complete source code and seed-controlled training scripts are
publicly available (Declarations: Code availability).

\subsection{Loss Function Formulation}
\label{sec:loss}

The total PINN loss function is the weighted sum of four terms:
\begin{equation}
  \mathcal{L}(\theta) = \lambda_r \mathcal{L}_r
  + \lambda_i \mathcal{L}_i
  + \lambda_b \mathcal{L}_b
  + \lambda_d \mathcal{L}_d,
  \label{eq:loss_total}
\end{equation}
The scalar weights $\lambda_r, \lambda_i, \lambda_b, \lambda_d \ge 0$
control the relative contribution of each term; imbalanced
weights induce gradient pathologies that prevent
convergence~\cite{Wang2021}. The individual terms are defined
as follows.

\paragraph{Residual loss $\mathcal{L}_r$} enforces the PDE
at $N_r$ collocation points
$\{(x_j^r, y_j^r, t_j^r)\}_{j=1}^{N_r}$:
\begin{equation}
  \mathcal{L}_r = \frac{1}{N_r} \sum_{j=1}^{N_r}
  \left| \frac{\partial \hat{C}}{\partial t}\bigg|_j
  - D\,\nabla^2 \hat{C}\bigg|_j
  + k\,\hat{C}_j \right|^2.
  \label{eq:loss_r}
\end{equation}

\paragraph{Initial condition loss $\mathcal{L}_i$} at
$N_i$ points:
\begin{equation}
  \mathcal{L}_i = \frac{1}{N_i} \sum_{j=1}^{N_i}
  \left| \hat{C}(x_j^i, y_j^i, 0) - C_0(x_j^i, y_j^i) \right|^2.
  \label{eq:loss_i}
\end{equation}

\paragraph{Boundary condition loss $\mathcal{L}_b$} at
$N_b$ boundary points:
\begin{equation}
  \mathcal{L}_b = \frac{1}{N_b} \sum_{j=1}^{N_b}
  \left| \mathcal{B}[\hat{C}](x_j^b, y_j^b, t_j^b) \right|^2,
  \label{eq:loss_b}
\end{equation}
where $\mathcal{B}[\cdot]$ denotes the boundary condition
operator corresponding to \eqref{eq:bc_left}--\eqref{eq:bc_right}.

\paragraph{Data loss $\mathcal{L}_d$} (for inverse problem;
set $\lambda_d = 0$ for forward problem) at $N_d$ observation
points:
\begin{equation}
  \mathcal{L}_d = \frac{1}{N_d} \sum_{j=1}^{N_d}
  \left| \hat{C}(x_j^d, y_j^d, t_j^d) - C_j^{\text{obs}} \right|^2.
  \label{eq:loss_d}
\end{equation}

We upweight the residual, using
$(\lambda_r, \lambda_i, \lambda_b) = (10, 1, 1)$ for the forward
problem ($\lambda_d = 0$) and $(10,1,1,1)$ for the inverse
problem, keeping the residual on the same gradient scale as the
other terms and preventing the optimizer from fitting
observations at the expense of physical consistency. Weights are
fixed throughout; adaptive schemes~\cite{Wang2021} are future
work.

\subsection{Collocation Point Sampling}
\label{sec:sampling}

Collocation points are sampled from $\Omega \times (0, T]$ using
Latin Hypercube Sampling, which provides better coverage than
uniform random sampling~\cite{McKay1979}: $N_r = 10{,}000$
interior points, $N_i = 400$ initial-condition points, and
$N_b = 400$ boundary points ($100$ per edge), plus $5{,}000$
held-out test points for monitoring generalization. All samples
regenerate from a fixed seed for reproducibility.

\subsection{Training Procedure}
\label{sec:training}

Training proceeds in two phases: Adam~\cite{Kingma2015} for
$20{,}000$ iterations at fixed learning rate $10^{-3}$, followed
by limited-memory BFGS~\cite{Liu1989} via
jaxopt~\cite{Blondel2022} to convergence (max.\ $5{,}000$
iterations, gradient tolerance $10^{-9}$), exploiting
second-order curvature once iterates enter a local minimum's
neighborhood. Loss weights follow Section~\ref{sec:loss} and
are held fixed throughout; adaptive reweighting
schemes~\cite{Wang2021} are left for future work.

All experiments use a fixed seed (1234) for floating-point
reproducibility. Forward-problem training completes in
$3$--$5$ minutes on a T4 GPU ($10$--$15$ on CPU); the full
reproduction pipeline (forward training, validation, five-seed
inverse ensemble, figure generation) completes in
$\sim 50$ minutes on a T4 GPU (Online Resource 1 for detail).

\subsection{Inverse Problem Formulation}
\label{sec:inverse}

The forward problem of Section~\ref{sec:math} assumes $D,k$ are
known and seeks $C(x,t)$; the inverse setting reverses this,
recovering $(D,k)$ from a partial, noisy observation. Classical
mesh-based solvers are expensive here, requiring a separate
nested forward solve per candidate $(D,k)$, with cost growing
linearly in candidates~\cite{Wiencke2020}. PINNs sidestep this by
treating $D$ and $k$ as additional trainable network parameters,
so one training run recovers both the field and the physics
constants~\cite{Raissi2019, Haghighat2021}.

Synthetic observations
$\{(x_j^d, y_j^d, t_j^d, C_j^{\text{obs}})\}_{j=1}^{N_d}$ are
generated by sampling locations uniformly from
$\Omega \times [t_{\min}, t_{\max}]$ with $t_{\min} = 0.5\,$ms
(just past the initial-condition transient, so data carry
information about both $D$ and $k$) and
$t_{\max} = 0.9\,T = 18\,$ms (the window over which the pulse
remains well-confined; cf.\ Online Resource 2). At each
location, the true concentration is computed by interpolating
the finite-difference reference solver (at $D=D^*$, $k=k^*$)
and corrupted by additive Gaussian noise,
\begin{equation}
  C_j^{\text{obs}} = C^*(x_j^d, y_j^d, t_j^d) + \varepsilon_j,
  \quad \varepsilon_j \sim \mathcal{N}\!\left(0,\,(\eta\, C_0)^2\right),
  \label{eq:observation}
\end{equation}
with $\eta$ the noise level as a fraction of $C_0$. Using the
finite-difference (rather than analytical) reference keeps the
observations exactly consistent with the Neumann boundary
conditions enforced by the PINN, avoiding a data-model mismatch
that would otherwise bias recovery at large $t$. Unless stated
otherwise, $N_d = 400$ and $\eta = 0.02$ ($2\%$, comparable to
high-quality FSCV noise floors~\cite{Rodeberg2017}).

$D$ and $k$ are introduced as trainable variables in the same
back-propagation graph as $\theta$, parametrized in log space
($D = \exp(\tilde{D})$, $k = \exp(\tilde{k})$) to enforce
positivity. Writing $\Phi = (\theta, \tilde{D}, \tilde{k})$, the
inverse PINN minimizes
\begin{equation}
  \Phi^* \;=\; \arg\min_{\Phi}\;
    \big[\lambda_r \mathcal{L}_r(\Phi)
       + \lambda_i \mathcal{L}_i(\Phi)
       + \lambda_b \mathcal{L}_b(\Phi)
       + \lambda_d \mathcal{L}_d(\Phi)\big],
  \label{eq:inverse_loss}
\end{equation}
with $(\lambda_r, \lambda_i, \lambda_b, \lambda_d) = (10,1,1,1)$,
so the optimizer prioritizes physical consistency over fitting
observation noise. Trainable parameters are initialized
perturbed from truth ($\exp(\tilde{D}_0) = 0.30$,
$\exp(\tilde{k}_0) = 0.04$); the optimization schedule matches
the forward problem (Section~\ref{sec:training}).

This is motivated by Section~\ref{sec:intro}: both parameters
are expected to shift in Parkinson's disease, and FSCV delivers
exactly the sparse, noisy measurements the inverse PINN consumes
\cite{Rodeberg2017}. This study restricts itself to controlled
synthetic validation; extending to real FSCV recordings is a
required next step, discussed in Section~\ref{sec:discussion}.

\section{Results}
\label{sec:results}

\subsection{Forward Problem: Accuracy and Convergence}
\label{sec:res_forward}

Fig~\ref{fig:forward} compares the trained PINN prediction
$\hat{C}(x, t; \theta)$ with the closed-form analytical
solution and the explicit finite-difference reference at three
representative time snapshots ($t = 0$, $T/4 = 5$\,ms,
$T/2 = 10$\,ms). The three curves coincide to within line
thickness across the entire spatial domain at every snapshot.
The full space-time error field (Fig~\ref{fig:heatmap}) shows the
pointwise absolute error $|C - \hat{C}|$ remains several orders
of magnitude below $C_0$ across the entire evaluation grid.

The PINN attains an L2 relative error of $1.52\%$ against a
finite-difference (FD) reference, integrated across the full
$201 \times 201 \times 200$ space-time grid, and $3.75\%$
against the closed-form infinite-domain analytical solution
(Online Resource 2) at $t=1$\,ms. The FD comparison is the
appropriate benchmark, since both solve the same bounded-domain
problem; it stays below $2.5\%$ at every reported time slice
(Table~\ref{tab:forward_error}). \emph{Errors against the
analytical reference grow substantially at later times}
($\sim 290\%$ at $t=20$\,ms), reflecting the expected breakdown
of the infinite-domain approximation as boundary reflections
(correctly modeled by the PINN) become non-negligible, not PINN
degradation. Training completes within $3$--$5$\,minutes on a
T4 GPU, with the loss following the two-phase Adam$\to$L-BFGS
pattern and no loss-weight pathology~\cite{Wang2021} under the
fixed $(10,1,1)$ schedule of Section~\ref{sec:training}.

\begin{figure}[t]
  \centering
  \includegraphics[width=0.82\textwidth]{02_Figures/Fig2_preview.png}
  \caption{\textbf{Forward-problem solution: PINN vs.\
  analytical and finite-difference references.}
  $C(x, t)$ at $t = 0$, $T/4 = 5$\,ms, and $T/2 = 10$\,ms.
  Analytical solution (solid black), finite-difference reference
  (dotted blue), and PINN prediction $\hat{C}$ (dashed red) are
  visually indistinguishable at all three slices. Parameters:
  $D = 0.32\,\mu\text{m}^2/\text{ms}$, $k = 0.05\,\text{ms}^{-1}$,
  $L = 5\,\mu\text{m}$, $\sigma = 0.5\,\mu\text{m}$,
  $C_0 = 1\,\mu\text{M}$}
  \label{fig:forward}
\end{figure}

\begin{figure}[t]
  \centering
  \includegraphics[width=0.82\textwidth]{02_Figures/Fig3_preview.png}
  \caption{\textbf{Space-time concentration field for the
  forward problem.} \textbf{Left}: analytical reference solution
  $C(x,t)$ on $(x,t) \in [-L/2,L/2] \times [0,T]$.
  \textbf{Centre}: PINN prediction $\hat{C}(x,t;\theta)$ on the
  same grid. \textbf{Right}: pointwise absolute error
  $|C - \hat{C}|$, colour-scaled independently. The release
  pulse spreads diffusively while decaying exponentially under
  the reuptake sink; the PINN reproduces both effects with
  errors several orders of magnitude below $C_0$}
  \label{fig:heatmap}
\end{figure}

\begin{table}[t]
\centering
\caption{\textbf{Forward-problem accuracy at representative
time snapshots.}
Relative $L^2$ error of $\hat{C}(x,y,t)$ vs.\ the FD reference
(meaningful benchmark, both solve the same bounded-domain
problem) and vs.\ the infinite-domain analytical solution (shown
for completeness; its error grows with $t$ as boundary
reflections, absent from that formula, become non-negligible).}
\label{tab:forward_error}
\begin{tabular}{lcc}
\toprule
\textbf{Time (ms)} & \textbf{PINN vs.\ FD reference}
                    & \textbf{PINN vs.\ analytical (infinite-domain)} \\
\midrule
$t = 1$   & $2.43\%$    & $3.75\%$   \\
$t = 5$   & $1.87\%$    & $51.97\%$  \\
$t = 10$  & $1.57\%$    & $129.26\%$ \\
$t = 20$  & $1.65\%$    & $290.43\%$ \\
\botrule
\end{tabular}
\end{table}

\subsection{Inverse Problem: Parameter Recovery}
\label{sec:res_inverse}

At the canonical operating point of $\eta = 0.02$ ($2\%$
Gaussian noise on $N_d = 400$ synthetic observations), the PINN
recovers $D = 0.3515 \pm 0.0034\,\mu\text{m}^2/\text{ms}$ and
$k = 0.0573 \pm 0.0006\,\text{ms}^{-1}$ against
$D^* = 0.32$, $k^* = 0.05$ (mean relative errors
$9.85\% \pm 1.07\%$, $14.61\% \pm 1.13\%$, five-seed ensemble;
Table~\ref{tab:inverse}). Values lie within canonical
experimental ranges for striatal
dopamine~\cite{CraggRice2004, Nicholson1981}, and the small
inter-seed SD ($\sim 1\%$) confirms the recovery is stable to
initialization rather than a lucky run. Single-seed variants reported elsewhere
(Tables~\ref{tab:noise_sweep}/\ref{tab:param_grid},
$9.23$--$9.29\%$ for $D$) sit comfortably within this ensemble's
$\pm 1.07\%$ envelope.

\begin{table}[t]
\centering
\caption{\textbf{Inverse-problem recovery accuracy at the
canonical operating point, averaged across five
network-initialization seeds.}
Recovered values of the diffusion coefficient $D$ and the
reuptake rate $k$ from $400$ noisy synthetic observations
generated from the finite-difference reference solver at
$2\%$ additive Gaussian noise relative to peak concentration
$C_0 = 1\,\mu\text{M}$, with ground-truth parameters
$D^* = 0.32\,\mu\text{m}^2/\text{ms}$ and
$k^* = 0.05\,\text{ms}^{-1}$. Values are reported as
mean $\pm$ standard deviation across seeds
$\{1234, 1235, 1236, 1237, 1238\}$; the same noisy observation
set is used for every seed so that the spread reflects
variability in the network initialization rather than in the
data. Both parameters are recovered within $15$ percent of
the truth. Relative error is computed against the true
parameter value.}
\label{tab:inverse}
\begin{tabular}{lcccc}
\toprule
\textbf{Noise level} & \textbf{Recovered $D$} & \textbf{Rel.\ error $D$}
& \textbf{Recovered $k$} & \textbf{Rel.\ error $k$} \\
\midrule
$2\%$    & $0.3515 \pm 0.0034$ & $9.85\% \pm 1.07\%$
         & $0.0573 \pm 0.0006$ & $14.61\% \pm 1.13\%$ \\
\botrule
\end{tabular}
\end{table}

Fig~\ref{fig:inverse} shows both parameters starting from
offset initial guesses and following monotone trajectories
toward truth during Adam. $D$, which couples to the second
spatial derivative of $C$ and so receives a strong gradient
signal from every collocation point, converges within the first
few thousand iterations; $k$ couples only to the concentration
itself and converges more gradually, with L-BFGS providing minor
refinement within the basin located by Adam.

This noise-robustness is what makes the framework practically
deployable: individual-subject dopamine measurements are
inevitably noisy, and a pipeline sensitive only to the order of
magnitude of measurement noise (rather than its precise value)
translates more easily across heterogeneous experimental
contexts than one tuned to a specific noise budget.

\begin{figure}[t]
  \centering
  \includegraphics[width=0.85\textwidth]{02_Figures/Fig4_preview.png}
  \caption{\textbf{Inverse-problem parameter recovery
  (primary-seed trajectory).}
  Trajectories of $D$ (\textbf{left}, red) and $k$
  (\textbf{right}, blue) vs.\ training iteration, from $400$
  noisy observations at $2\%$ noise. Dotted lines mark ground
  truth $D^* = 0.32$, $k^* = 0.05$. Both start offset
  ($D_0 = 0.30$, $k_0 = 0.04$) and approach truth during Adam,
  refined under L-BFGS, settling slightly above truth (the
  diffusion-reuptake coupling of Section~\ref{sec:lim_coupling});
  the other four seeds produce superimposable trajectories (not
  shown)}
  \label{fig:inverse}
\end{figure}

\subsection{Noise Sensitivity}
\label{sec:res_noise}

We re-ran the inverse PINN at five noise levels
$\eta \in \{0\%, 1\%, 2\%, 5\%, 10\%\}$ relative to $C_0$, using
the same $400$ observation locations and seed at every level
(Table~\ref{tab:noise_sweep}, Fig~\ref{fig:noise_sweep}).
Recovery is essentially exact at $\eta = 0\%$ ($1.45\%$ error
for $D$, $1.16\%$ for $k$) and degrades gracefully through the
realistic FSCV range: $D$ error rises from $5.5\%$ ($\eta=1\%$)
to $9.3\%$ ($\eta=2\%$) to $30.4\%$ ($\eta=5\%$), while $k$
error rises from $8.3\%$ to $14.9\%$ before plateauing near
$13.5\%$. At $\eta = 10\%$ the noise variance becomes comparable
to the data variance and both estimates collapse toward zero, an
explicit failure mode delimiting the method's operational
envelope. The non-monotonic $k$ error between $\eta=2\%$ and
$\eta=5\%$ reflects the diffusion-reuptake coupling characterized
in Section~\ref{sec:lim_coupling}. For realistic FSCV recordings,
where instrument noise floors are typically $1$--$3\%$
\cite{Rodeberg2017}, the inverse PINN remains within the $15\%$
accuracy envelope reported as the canonical result.

\begin{table}[t]
\centering
\caption{\textbf{Noise sensitivity of the inverse-problem
recovery.} Recovered $D$ and $k$ from $400$ noisy synthetic
observations at five noise levels $\eta$ (percentages of $C_0$),
primary seed ($1234$).}
\label{tab:noise_sweep}
\begin{tabular}{lcccc}
\toprule
\textbf{Noise} & \textbf{Recovered $D$} & \textbf{Rel.\ error $D$}
& \textbf{Recovered $k$} & \textbf{Rel.\ error $k$} \\
\midrule
$0\%$ (clean) & $0.325$ & $1.45\%$  & $0.0506$ & $1.16\%$  \\
$1\%$         & $0.338$ & $5.50\%$  & $0.0542$ & $8.31\%$  \\
$2\%$         & $0.350$ & $9.29\%$  & $0.0574$ & $14.85\%$ \\
$5\%$         & $0.417$ & $30.40\%$ & $0.0568$ & $13.52\%$ \\
$10\%$        & $0.195$ & $39.17\%$ & $0.0125$ & $75.06\%$ \\
\botrule
\end{tabular}
\end{table}

\begin{figure}[t]
  \centering
  \includegraphics[width=0.8\textwidth]{02_Figures/Fig5_preview.png}
  \caption{\textbf{Noise sensitivity of inverse-problem
  recovery.} Recovered $D$ (left) and $k$ (right) vs.\
  observational noise level $\eta$. Dotted lines mark ground
  truth. Recovery is essentially exact at $\eta = 0\%$ and
  degrades gracefully through the realistic FSCV regime
  ($\eta \sim 2\%$) and the moderate condition ($\eta = 5\%$);
  at $\eta = 10\%$ both estimates collapse toward zero}
  \label{fig:noise_sweep}
\end{figure}

\subsection{Bayesian Uncertainty Quantification}
\label{sec:res_uq}

The multi-seed analysis of Section~\ref{sec:res_inverse}
quantifies initialization-to-initialization variability; to
additionally quantify the \emph{data-driven} uncertainty from
the $400$ observations, we applied a Laplace approximation
around the MAP estimate and, as a cross-check, an HMC
sampler~\cite{cabezas2024blackjax}, both using the
infinite-domain analytical solution (Online Resource 2) as the
likelihood forward model. The two agree to three decimal places
(Online Resource 3); we report Laplace here.

\begin{figure}[t]
  \centering
  \includegraphics[width=0.7\textwidth]{02_Figures/Fig6_preview.png}
  \caption{\textbf{Bayesian posterior over $(D, k)$
  (Laplace approximation, analytical-model likelihood).}
  Blue: $2{,}000$ Laplace posterior samples. Navy ellipses:
  $1$--$3\sigma$ contours. Red star: analytical MAP. Green
  square: PINN point estimate. Black plus: ground truth
  $(D^*, k^*) = (0.32, 0.05)$}
  \label{fig:posterior}
\end{figure}

The Laplace approximation yields
$D = 0.320 \pm 0.011\,\mu\text{m}^2/\text{ms}$, $95\%$ CI
$[0.298, 0.343]$, containing $D^* = 0.32$ -- usable directly as
a data-implied uncertainty estimate. For $k$, both methods
collapse to the boundary $k \to 0$, mirroring the unphysical
$k_{\mathrm{LM}} = -0.040$ from the LM baseline
(Section~\ref{sec:res_lm_baseline}): the infinite-domain model
cannot reproduce elevated late-time concentrations in the
bounded-domain data, so maximum-likelihood $k$ is pushed toward
zero, whereas the PINN's bounded-domain residual loss returns
physically meaningful $k$. We report Bayesian intervals for $D$
only and rely on the multi-seed ensemble for $k$; a fully
consistent Bayesian $k$ uncertainty would require
profile-likelihood PINN retraining at each $(D,k)$, left as
future work.

\subsection{Parameter Recovery Across the Biological Range}
\label{sec:res_param_grid}

To verify that inverse-problem accuracy is not specific to the
canonical operating point, we re-ran the inverse PINN at five
additional $(D^*, k^*)$ combinations spanning the
healthy-to-pathological striatal range
\cite{CraggRice2004, Nicholson1981, Wiencke2020}
($D \in \{0.16, 0.32, 0.50\}\,\mu\text{m}^2/\text{ms}$,
$k \in \{0.025, 0.05, 0.10\}\,\text{ms}^{-1}$;
Table~\ref{tab:param_grid}, Fig~\ref{fig:param_grid}). The mean
recovery error is $7.9\%$ for $D$ and $18.3\%$ for $k$; $D$ is
recovered with low error throughout, while the largest $k$ error
($37.1\%$ at $D^*=0.50$, $k^*=0.025$) occurs when the reuptake
timescale $1/k^*$ exceeds the simulation window $T$, so the data
carries weak information about $k$ and the optimizer compensates
by adjusting $D$ along the trade ray
(Section~\ref{sec:lim_coupling}). Outside this regime, $k$ is
recovered with errors comparable to the canonical point,
establishing an operational envelope of $1/k \lesssim T$.

\begin{table}[t]
\centering
\caption{\textbf{Parameter-recovery accuracy across the
biological range.} $400$ noisy observations at $2\%$ noise,
seed $1234$.}
\label{tab:param_grid}
\begin{tabular}{cccccc}
\toprule
\textbf{$D^*$} & \textbf{$k^*$}
& \textbf{$D_{\mathrm{rec}}$} & \textbf{$|$err$|_D$ (\%)}
& \textbf{$k_{\mathrm{rec}}$} & \textbf{$|$err$|_k$ (\%)} \\
\midrule
$0.16$ & $0.025$ & $0.1624$ & $1.49$  & $0.0295$ & $18.14$ \\
$0.16$ & $0.10$  & $0.1661$ & $3.84$  & $0.1068$ & $6.75$  \\
$0.32$ & $0.05$  & $0.3495$ & $9.23$  & $0.0571$ & $14.17$ \\
$0.50$ & $0.025$ & $0.5592$ & $11.83$ & $0.0343$ & $37.12$ \\
$0.50$ & $0.10$  & $0.5655$ & $13.09$ & $0.1150$ & $15.05$ \\
\botrule
\end{tabular}
\end{table}

\begin{figure}[t]
  \centering
  \includegraphics[width=0.8\textwidth]{02_Figures/Fig7_preview.png}
  \caption{\textbf{Parameter recovery across the biological
  range.} Recovered $D$ (left) and $k$ (right) vs.\ ground truth
  across the five $(D^*,k^*)$ combinations above. Dashed line is
  $y=x$ (perfect recovery)}
  \label{fig:param_grid}
\end{figure}

\subsection{Observation Density Scaling}
\label{sec:res_obs_density}

We separately re-ran the inverse PINN at observation counts
$N_{\mathrm{obs}} \in \{100, 200, 400, 800, 1600\}$ at the
canonical operating point (Fig~\ref{fig:obs_density}). Going
from $N_{\mathrm{obs}}=100$ to $1600$ reduces mean error from
$24.8\%$ to $7.0\%$ for $D$ and $24.7\%$ to $4.7\%$ for $k$,
broadly consistent with $N_{\mathrm{obs}}^{-1/2}$ Cram\'er--Rao
scaling, before plateauing at a coupling-imposed bias floor of
$\sim 5\%$ for both parameters, the empirical signature of the
diffusion-reuptake identifiability limit: more measurements do
not break the trade ray. Field investigators should expect a
$\sim 10$--$15\%$ irreducible parameter uncertainty regardless
of measurement quantity, absent auxiliary constraints such as
observed peak time or amplitude.

\begin{figure}[t]
  \centering
  \includegraphics[width=0.85\textwidth]{02_Figures/Fig8_preview.png}
  \caption{\textbf{Observation density scaling.} Relative
  recovery error for $D$ (red) and $k$ (blue) vs.\ number of
  observations $N_{\mathrm{obs}}$, log-log axes, with dotted
  $N_{\mathrm{obs}}^{-1/2}$ Cram\'er--Rao reference lines. Error
  follows the asymptotic scaling at low-to-moderate
  $N_{\mathrm{obs}}$, then plateaus at a coupling-imposed bias
  floor}
  \label{fig:obs_density}
\end{figure}

\subsection{Comparison with a Levenberg--Marquardt Baseline}
\label{sec:res_lm_baseline}

The natural classical baseline is to fit the closed-form
analytical solution (infinite-domain 2D Gaussian) to the noisy
observations via Levenberg--Marquardt (LM) non-linear
least-squares, using \texttt{scipy.optimize.least\_squares} on
the same canonical $400$-observation dataset
(Table~\ref{tab:lm_baseline}), as a control clarifying what the
PINN specifically contributes.

\begin{table}[t]
\centering
\caption{\textbf{Comparison of inverse-PINN recovery against the
Levenberg--Marquardt baseline.} LM fits the infinite-domain
analytical Gaussian; the PINN solves the bounded-domain PDE
without assuming an analytical form.}
\label{tab:lm_baseline}
\begin{tabular}{lcc}
\toprule
                                   & \textbf{LM (analytical)} & \textbf{PINN (numerical)} \\
\midrule
$D$ recovered                      & $0.4318$                 & $0.3495$                  \\
$D$ rel.\ error (\%)               & $34.92$                  & $9.23$                    \\
$k$ recovered                      & $-0.0403$ (unphysical)   & $0.0571$                  \\
$k$ rel.\ error (\%)               & $180.6$ (unphysical)     & $14.17$                   \\
Runtime                            & $\sim 3.3$\,ms           & $\sim 540$\,s             \\
Handles arbitrary domain geometry  & no                       & yes (FD-validated)        \\
\botrule
\end{tabular}
\end{table}

Although the LM fit converges in $\sim 3$ milliseconds (five
orders of magnitude faster), it recovers a significantly biased
$D$ ($34.9\%$ error vs.\ the PINN's $9.2\%$) and a
\emph{negative} reuptake rate ($k_{\mathrm{LM}} =
-0.040\,\text{ms}^{-1}$) -- a physically impossible value. This
failure is structural: the infinite-domain solution assumes free
spreading, whereas in the bounded domain $[-L/2, L/2]^2$
dopamine accumulates at late times because it cannot exit, so LM
compensates by driving $k$ negative. The PINN avoids this by
learning the actual bounded-domain dynamics rather than relying
on an inappropriate analytical surrogate; its single unified
surrogate for the forward and inverse problems, extending
naturally to non-analytical settings, is the principal practical
contribution of this work.

\subsection{Architecture Ablation}
\label{sec:res_arch_ablation}

Re-training the forward problem with three MLP configurations
(compact, canonical, wide) confirms the canonical four-layer,
64-neuron architecture is not badly over- or under-parameterized:
all three achieve sub-$2.1\%$ forward $L^2$ error against the FD
reference, with the wide configuration slightly less accurate
($2.04\%$ vs.\ $1.52\%$), likely reflecting mild
overparameterization. Full configuration table in Online
Resource 4E.

\section{Discussion}
\label{sec:discussion}

\subsection{Biological Interpretation of the Recovered Parameters}
\label{sec:disc_bio}

The recovered $D = 0.3515 \pm 0.0034\,\mu\text{m}^2/\text{ms}$
and $k = 0.0573 \pm 0.0006\,\text{ms}^{-1}$ fall within canonical
experimental ranges for striatal
dopamine~\cite{CraggRice2004, Nicholson1981, Rice2008},
indicating physiologically meaningful recovery rather than
overfitting. Both are expected to shift in Parkinson's disease
 $D$ via tortuosity changes, $k$ via DAT loss underlying
DAT-PET imaging~\cite{Rice2008, CraggRice2004}, but this study
does not measure disease-associated change; it demonstrates only
that the method can recover $D$ and $k$ from the sparse, noisy
regime FSCV produces. Whether the constants can distinguish
healthy from Parkinsonian signalling on real recordings is
follow-up work (Section~\ref{sec:disc_future}); the contribution
here is the inverse formulation, not a validated biomarker.

\subsection{Comparison with Prior Work}
\label{sec:disc_priorwork}

Compared with finite-difference/finite-element approaches to
synaptic dopamine modeling~\cite{Wiencke2020}, the PINN
framework eliminates mesh-construction overhead and unifies
forward and inverse problems in a single optimization. Within
the broader landscape of biomedical PINN
applications~\cite{Sahli2020, Yazdani2020, Haghighat2021}, this
study is, to our knowledge, an early example targeting
neurotransmitter diffusion at the synaptic scale.

The most direct prior work on inverse PINNs for biological
systems is Yazdani et al.~\cite{Yazdani2020}, targeting lumped
ODE reaction kinetics in well-mixed compartments, whereas the
present work tackles a spatiotemporal PDE where
$D\,\nabla^2 C$ couples directly to second spatial derivatives.
Yazdani et al.\ report $5$--$30\%$ errors for identifiable
parameters, with others practically non-identifiable via Fisher
Information Matrix (FIM) null-eigenvector analysis; our errors
of $9.85\% \pm 1.07\%$ ($D$) and $14.61\% \pm 1.13\%$ ($k$) lie
in the same band while addressing a genuinely spatial problem.
The diffusion-reuptake coupling biasing our recovery
(Section~\ref{sec:lim_coupling}) is the spatial analogue of
their FIM-null-eigenvector limits, suggesting practical
identifiability recurs across both lumped and spatially-resolved
inverse PINNs.

\subsection{Inverse-Problem Identifiability and Parameter Conditioning}
\label{sec:identifiability}

Independently of the solution method (PINN, non-linear
least-squares, Bayesian inversion), the well-posedness of the
$(D,k)$ inverse problem for equation~\eqref{eq:pde} is governed
by \emph{structural identifiability} (do distinct parameters
produce distinguishable fields?) and \emph{practical
identifiability} (do noisy, finite data localize them usefully?)
\cite{Yazdani2020}. The forward map is structurally identifiable
on dense, noise-free data; practical identifiability is more
subtle and is the controlling factor for the recovery accuracy
reported throughout Section~\ref{sec:results}.

Near $(D^*, k^*)$, the concentration response to leading order is
\begin{equation*}
  \Delta C \;\approx\; \frac{\partial C}{\partial D} \Delta D
                     + \frac{\partial C}{\partial k} \Delta k,
\end{equation*}
and $\partial C/\partial D$, $\partial C/\partial k$ are not
linearly independent over the observable space-time region: both
diffusion and reuptake reduce concentration at a fixed point, so
a small increase in $D$ paired with a small increase in $k$
perturbs $C$ by less than the $2\%$ noise floor within the
observation window $t \in [0.5, 18]\,\text{ms}$. The inverse
problem thus exhibits a one-dimensional \emph{trade-off
manifold},  the \emph{diffusion-reuptake trade ray}  along
which the loss is flat to within the noise scale.

This ray is structural, not a finite-sample artefact: a
Cram\'er--Rao analysis predicts error decreasing as
$N_{\mathrm{obs}}^{-1/2}$ only when the Fisher information
matrix (FIM) is non-singular, but here the FIM has a small
eigenvalue in the trade-ray direction, so more observations
sharpen the well-constrained direction while barely improving
the other, confirmed by the observation-density sweep
(Section~\ref{sec:res_obs_density}), which asymptotes to a
nonzero bias floor. Tighter recovery needs qualitatively
different measurements ($t \gg 1/k$, peak amplitude/time, or
drug-blockade conditions), not more of the same. The noise sweep
(Section~\ref{sec:res_noise}) adds a complementary view: recovery
holds within $15\%$ at $\eta = 2\%$ but collapses by
$\eta = 10\%$; and the parameter grid shows that when
$1/k^* > T$, the data carries too little information about $k$
regardless of noise level.

\paragraph{Empirical signature of the coupling.}
\label{sec:lim_coupling}
The trade-ray structure is directly visible in the multi-seed
ensemble (Section~\ref{sec:res_inverse}): across all five seeds,
$D$ and $k$ are simultaneously biased \emph{above} truth
($9.9\% \pm 1.1\%$ and $14.6\% \pm 1.1\%$) in matched
directions, with inter-seed SD under $1.1\%$ -- a systematic
feature of the loss landscape, not optimizer drift. The Bayesian
analysis (Section~\ref{sec:res_uq}) corroborates this: the
Laplace posterior collapses $k$ to the boundary $k \to 0$ while
leaving $D$ well-constrained, mirroring the unphysical
$k_{\mathrm{LM}} = -0.04$ from the Levenberg--Marquardt baseline
(Section~\ref{sec:res_lm_baseline}), two independent
manifestations of the same identifiability limit.

\paragraph{Implications for biological interpretation.}
Recovered $(D, k)$ should be interpreted \emph{jointly} rather
than as independent estimates; comparisons across conditions
(e.g.\ healthy vs.\ Parkinsonian signalling) should examine the
joint direction of change along the trade ray rather than
marginal $D$ or $k$ alone; and experimental protocols that break
the trade ray, extended observation windows, paired
iontophoretic release, or simultaneous receptor-binding
measurements are the most effective way to improve recovery
beyond the bias floor characterised here.

\subsection{Limitations and Scope}
\label{sec:limitations}

We explicitly delineate the scope and the limitations of the
present study.

\paragraph{Reduced spatial dimensionality.}
The model is restricted to a planar 2D domain and a
single-pulse release event; real perisynaptic geometry is fully
3D with structural heterogeneity this abstraction omits. The 2D
restriction follows prior modeling
literature~\cite{Wiencke2020} and keeps cost tractable;
extension to 3D is a straightforward modification of the input
layer and PDE residual, but whether the reported accuracy
persists in the heterogeneous 3D regime remains to be shown.

\paragraph{Synthetic observation data and absence of FSCV
validation.}
The inverse problem is validated only against synthetic noisy
observations from a finite-difference reference of the same
governing PDE, which makes recovery easier than from genuine
experimental data subject to instrument response and unmodeled
biological variability. \emph{No validation against real FSCV
recordings has been performed}, and we make no claim the
reported accuracy transfers to in vivo data; such validation,
on data from healthy and Parkinsonian
subjects~\cite{Rodeberg2017}, is a required next step
(Section~\ref{sec:disc_future}).

\paragraph{Linearised kinetics, homogeneous medium, no receptor
binding.}
Reuptake is first-order, whereas DAT transport is more
accurately Michaelis--Menten with saturation~\cite{Wiencke2020,
CraggRice2004}; this holds below $K_m \approx 0.2\,\mu$M but can
underestimate clearance immediately after release. $D$ is a
single effective scalar despite real cell-scale medium
heterogeneity~\cite{CraggRice2004, Nicholson1981}, and the model
excludes postsynaptic D1/D2 receptor binding, consistent with
how the dopamine literature typically separates it from
extracellular diffusion~\cite{Rice2008}. Extensions to
Michaelis--Menten kinetics or spatially-varying $D(x,y)$ are
methodologically straightforward but would each aggravate the
identifiability issue above; we leave them to follow-up work.

\paragraph{Single operating point.}
Results use a single hyperparameter setting; seed variability is
small (Table~\ref{tab:inverse}, $\sigma_D/D, \sigma_k/k < 1\%$),
but a systematic sweep over collocation density, network depth,
and adaptive loss-weight schemes~\cite{Wang2021} would further
strengthen robustness.

\subsection{Future Directions}
\label{sec:disc_future}

These limitations delineate the most immediate paths for future
work: extending to fully 3D synaptic geometries from electron
microscopy, to test whether 2D forward accuracy persists in the
heterogeneous 3D regime; and, most importantly, replacing
synthetic data with real FSCV measurements~\cite{Rodeberg2017}, 
a prerequisite for clinical interpretation, for which the
inverse PINN's data-efficiency is well suited but as yet
undemonstrated. Longer term, coupling this local model with
network-level descriptions of alpha-synuclein spread and
dopaminergic neuron loss~\cite{Garcia2017, Corti2023,
Pandya2019} would support a multi-scale pipeline linking
molecular signalling deficits to clinical decline.

\section{Conclusion}
\label{sec:conclusion}

We presented and validated a physics-informed neural network
framework for the forward and inverse problems of dopamine
diffusion in the synaptic cleft. The forward PINN agrees with
the finite-difference reference across the full simulation
window ($1.52\%$ L2 error) and the analytical solution at early
times ($3.75\%$); the inverse formulation recovers $D$ and $k$
with mean relative errors of $9.85\% \pm 1.07\%$ and
$14.61\% \pm 1.13\%$ at $2\%$ observational noise. These results
establish PINNs as a mesh-free, data-efficient alternative to
classical solvers for synaptic neurotransmitter dynamics, with
the inverse formulation offering a tractable route to
subject-specific parameter estimation that avoids the costly
nested-optimization loops of mesh-based
pipelines~\cite{Wiencke2020}. Future work will extend to 3D
synaptic geometries, incorporate real FSCV
recordings~\cite{Rodeberg2017}, and couple this local model with
network-level descriptions of Parkinson's disease
progression~\cite{Garcia2017, Corti2023, Pandya2019}.

\section*{Declarations}

\bmhead{Funding}
The authors received no specific funding for this work.

\bmhead{Competing interests}
The authors have declared that no competing interests exist.

\bmhead{Ethics approval}
Not applicable. This study involved no human participants,
human data, or animals; all results derive from synthetic
observations generated by a finite-difference reference solver.

\bmhead{Consent to participate / Consent to publish}
Not applicable.

\bmhead{Data availability}
All data are publicly available. Synthetic datasets, trained
PINN model configurations, and analysis notebooks are deposited
in a public GitHub repository
(\url{https://github.com/Hachero98/dopamine-PINN-}) and archived
on Zenodo (DOI
\href{https://doi.org/10.5281/zenodo.20352595}{10.5281/zenodo.20352595}),
including a complete reference run with all figures and
\texttt{metrics.json}.

\bmhead{Code availability}
All Python (JAX/Flax NNX) code used to produce the results
reported in this manuscript is released under the MIT License
in the same repository referenced in the Data availability
statement. The repository includes a \texttt{README} with
reproduction instructions, a pinned \texttt{requirements.txt},
and seed-controlled scripts and a Colab notebook that
regenerate every figure and table in the paper.

\bmhead{Author contributions}
Conceptualization: E.H., H.Z.; Methodology: E.H.; Software: E.H.;
Formal analysis: E.H.; Investigation: E.H.; Writing - original
draft: E.H.; Writing -review and editing: E.H., H.Z.;
Supervision: H.Z. Both authors read and approved the final
manuscript.


\setlength{\bibsep}{0pt plus 0.3ex}
\bibliographystyle{sn-basic}
\bibliography{refs}

\end{document}
